% =========================================================
\subsubsection{Limit Points}
\label{sec:limit-points}
% =========================================================

% ---- Toolkit ----
\begin{tcolorbox}[colback=gray!6, colframe=gray!40, arc=2pt,
  left=6pt, right=6pt, top=4pt, bottom=4pt,
  title={\small\textbf{Limit Points --- Quick Reference}},
  fonttitle=\small\bfseries]
\small
\begin{tabular}{l l l l}
\toprule
\textbf{Label} & \textbf{Statement} & \textbf{Proof method} & \textbf{Detail} \\
\midrule
D\ref{def:limit-point}    & Limit point of $E$: every $B_r(x)$ meets $E\setminus\{x\}$        & Definition             & \hyperref[def:limit-point]{$\downarrow$} \\
D\ref{def:derived-set}    & $E' :=$ set of all limit points of $E$                            & Definition             & \hyperref[def:derived-set]{$\downarrow$} \\
D\ref{def:isolated-point} & $a\in E$ isolated iff $a\notin E'$                                & Definition             & \hyperref[def:isolated-point]{$\downarrow$} \\
T\ref{thm:seq-char-limit} & $x\in E' \iff \exists(x_n)\subseteq E\setminus\{x\}$, $x_n\to x$ & Sequential, $\varepsilon$-$\delta$ & \hyperref[thm:seq-char-limit]{$\downarrow$} \\
C\ref{cor:inf-nbhd}       & $x\in E' \iff$ every $B_r(x)\cap E$ is infinite                  & Seq.\ char.\ + construction & \hyperref[cor:inf-nbhd]{$\downarrow$} \\
T\ref{thm:derived-closed} & $E'$ is closed                                                    & Sequential closure     & \hyperref[thm:derived-closed]{$\downarrow$} \\
T\ref{thm:closed-derived} & $E$ closed $\iff E'\subseteq E$                                   & Sequential closure     & \hyperref[thm:closed-derived]{$\downarrow$} \\
T\ref{thm:closure-formula}& $\overline{E} = E\cup E'$                                         & Definition of closure  & \hyperref[thm:closure-formula]{$\downarrow$} \\
D\ref{def:perfect-set}    & $E$ perfect iff $E$ closed and $E=E'$                             & Definition             & \hyperref[def:perfect-set]{$\downarrow$} \\
\bottomrule
\end{tabular}
\end{tcolorbox}

\medskip
\begin{remark*}[Terminology note]
The term \emph{limit point} is standard in Rudin and Abbott.
Pugh uses \emph{cluster point} and \emph{accumulation point} as synonyms.
Be aware that some authors (including Pugh) use ``limit of $S$'' more broadly,
allowing the limit sequence to be constant; in that usage a finite set can have
limits.
The definition here follows Rudin: limit points require a \emph{distinct}
sequence, so a finite set has no limit points.
\end{remark*}

% ---- Limit Point ----
\begin{tcolorbox}[colback=propbox, colframe=propborder, arc=2pt,
  left=6pt, right=6pt, top=4pt, bottom=4pt,
  title={\small\textbf{Definition (Limit Point)}},
  fonttitle=\small\bfseries]
\begin{definition}[Limit Point]
\label{def:limit-point}
Let $(X,d)$ be a metric space and let $E \subseteq X$.
A point $x \in X$ is called a \emph{limit point} of $E$ if
\[
\forall r > 0,\quad B_r(x) \cap (E \setminus \{x\}) \neq \emptyset.
\]
\end{definition}
\end{tcolorbox}

\begin{remark*}[Fully quantified form]
\[
\forall (X,d)\ \text{metric space},\
\forall E \subseteq X,\
\forall x \in X,
\]
\[
x \text{ is a limit point of } E
\iff
\forall r \in \mathbb{R}^+,\
\exists y \in E \setminus \{x\}
\text{ such that } d(x,y) < r.
\]
\end{remark*}

% =========================================================
% TikZ diagram: Limit Point
% Place immediately after the Limit Point definition box.
% Requires: \usetikzlibrary{backgrounds, decorations.pathmorphing}
% =========================================================

\begin{figure}[h]
\centering
\begin{tikzpicture}[
  every label/.style = {font=\small},
  dot/.style         = {circle, fill, inner sep=1.8pt},
  opendot/.style     = {circle, draw, fill=white, inner sep=1.8pt, line width=0.8pt}
]

% --- Ambient space label ---
\node[font=\small\itshape, anchor=north west] at (-4.2, 3.3) {$X$};

% --- Set E (irregular blob, filled) ---
\begin{scope}
  \fill[blue!10, draw=blue!50, line width=0.8pt,
        decorate,
        decoration={random steps, segment length=6pt, amplitude=4pt}]
    (0.6, 0.2)
    .. controls (1.8, 2.2) and (0.2, 3.0) .. (-1.0, 2.8)
    .. controls (-2.4, 2.6) and (-2.8, 1.2) .. (-2.0, 0.2)
    .. controls (-1.2,-0.8) and ( 0.0,-0.6) .. ( 0.6, 0.2)
    -- cycle;
\end{scope}
\node[font=\small\itshape, text=blue!60!black] at (-1.1, 1.4) {$E$};

% --- Balls B_r(x) centred at x = (1.6, 1.0) ---
% x sits just outside / on the fringe of E
\coordinate (x) at (1.6, 1.0);

\draw[gray!60, dashed, line width=0.6pt] (x) circle (2.2cm);
\draw[gray!60, dashed, line width=0.6pt] (x) circle (1.4cm);
\draw[gray!60, dashed, line width=0.6pt] (x) circle (0.7cm);

% Ball labels (on the right arc of each circle)
\node[font=\scriptsize, text=gray!70!black] at ($(x) + (1.556, 1.556)$) [right] {$B_{r_1}(x)$};
\node[font=\scriptsize, text=gray!70!black] at ($(x) + (0.990, 0.990)$) [right] {$B_{r_2}(x)$};
\node[font=\scriptsize, text=gray!70!black] at ($(x) + (0.495, 0.495)$) [right] {$B_{r_3}(x)$};

% --- Points of E inside each punctured ball ---
% One witness point per ball, each clearly in E \ {x}
\node[dot, fill=blue!70!black, label={[blue!70!black]above left:\scriptsize$y_1$}]
  at ($(x) + (-1.5, 1.0)$) {};   % inside B_{r1}, inside E

\node[dot, fill=blue!70!black, label={[blue!70!black]above:\scriptsize$y_2$}]
  at ($(x) + (-0.7, 0.9)$) {};   % inside B_{r2}, inside E

\node[dot, fill=blue!70!black, label={[blue!70!black]above left:\scriptsize$y_3$}]
  at ($(x) + (-0.35, 0.55)$) {}; % inside B_{r3}, inside E

% --- The limit point x itself (open dot — not required to be in E) ---
\node[opendot, label={[anchor=north west, font=\small]north west:$x$}] at (x) {};

% --- Caption line beneath ---
\end{tikzpicture}

\smallskip
{\small
For every radius $r > 0$, the punctured ball $B_r(x) \cap (E \setminus \{x\})$ is
non-empty: witness points $y_1, y_2, y_3 \in E \setminus \{x\}$ are shown for
three nested balls.
Hence $x \in E'$, even though $x \notin E$.
}
\caption{A limit point $x$ of the set $E$ in a metric space $(X,d)$.}
\label{fig:limit-point}
\end{figure}




















\begin{remark*}[Intuition]
A limit point of $E$ is a point that $E$ crowds around, even if $x$ itself
is not in $E$.
Every punctured neighbourhood of $x$ must meet $E$.
The exclusion of $x$ itself prevents isolated points from qualifying.
\end{remark*}

% ---- Derived Set ----
\begin{tcolorbox}[colback=propbox, colframe=propborder, arc=2pt,
  left=6pt, right=6pt, top=4pt, bottom=4pt,
  title={\small\textbf{Definition (Derived Set)}},
  fonttitle=\small\bfseries]
\begin{definition}[Derived Set]
\label{def:derived-set}
Let $(X,d)$ be a metric space and let $E \subseteq X$.
The set of all limit points of $E$ is called the \emph{derived set} of $E$
and is denoted $E'$:
\[
E' := \{x \in X : x \text{ is a limit point of } E \}.
\]
\end{definition}
\end{tcolorbox}

\begin{remark*}[Fully quantified form]
\[
\forall (X,d)\ \text{metric space},\
\forall E \subseteq X,
\]
\[
E' = \{\, x \in X :
\forall r \in \mathbb{R}^+,\
\exists y \in E \setminus \{x\}\ \text{such that}\ d(x,y) < r
\,\}.
\]
\end{remark*}

% =========================================================
% TikZ diagram: Derived Set
% Place immediately after the Derived Set definition box.
% Requires in preamble: \usepackage{tikz}
%                       \usetikzlibrary{calc}
% =========================================================

\begin{figure}[h]
\centering
\begin{tikzpicture}[
  dot/.style     = {circle, fill, inner sep=1.8pt},
  opendot/.style = {circle, draw, fill=white, inner sep=1.8pt, line width=0.8pt},
  lbl/.style     = {font=\small},
  sublbl/.style  = {font=\scriptsize, text=gray!55!black}
]

% --- E (blue blob) ---
\fill[blue!12, draw=blue!50, line width=0.9pt]
  (-0.1, -0.5)
  .. controls ( 0.6, -0.8) and ( 1.4,  0.0) .. ( 1.3,  1.0)
  .. controls ( 1.2,  2.0) and ( 0.8,  3.0) .. (-0.2,  3.3)
  .. controls (-1.0,  3.6) and (-2.6,  3.4) .. (-3.2,  2.4)
  .. controls (-3.8,  1.4) and (-3.6,  0.2) .. (-2.8, -0.4)
  .. controls (-2.0, -1.0) and (-0.8, -0.8) .. (-0.1, -0.5)
  -- cycle;

\node[lbl, text=blue!60!black] at (-1.6, 0.5) {$E$};

% --- E' (orange dashed contour) ---
% Right lune (x > 0.9) lies in E' but not E  → E' \ E
% Top-left pocket (x < -2.5, y > 2.2) lies in E but not E' → E \ E'
\draw[orange!75!black, dashed, line width=1.2pt]
  ( 1.3,  0.0)
  .. controls ( 2.0,  0.6) and ( 2.0,  2.6) .. ( 1.1,  3.5)
  .. controls ( 0.2,  4.2) and (-1.4,  3.8) .. (-2.2,  3.0)
  .. controls (-2.8,  2.4) and (-2.6,  1.0) .. (-2.2,  0.2)
  .. controls (-1.8, -0.6) and ( 0.2, -0.8) .. ( 1.3,  0.0)
  -- cycle;

\node[lbl, text=orange!75!black] at (2.55, 1.8) {$E'$};

% --- Ambient space ---
\node[lbl, font=\small\itshape] at (-4.4, 3.8) {$X$};

% --- Region labels ---
\node[sublbl] at (-0.6, -0.1) {$E \cap E'$};
\node[sublbl, align=center] at ( 1.8,  0.6) {$E'{\setminus}E$};

% --- a in E ∩ E' : non-isolated interior point ---
\node[dot, fill=blue!70!black] (A) at (-1.0, 1.6) {};
\node[lbl, anchor=south west] at ($(A)+(0.10, 0.10)$) {$a$};
\node[sublbl, anchor=north west] at ($(A)+(0.10,-0.10)$)
  {$a\in E,\;a\in E'$};

% --- b in E' \ E : limit point of E not in E ---
\node[opendot, draw=orange!75!black] (B) at (1.05, 1.8) {};
\node[lbl, anchor=south west] at ($(B)+(0.10, 0.10)$) {$b$};
\node[sublbl, anchor=north west] at ($(B)+(0.10,-0.10)$)
  {$b\notin E,\;b\in E'$};

% --- c in E \ E' : isolated point of E ---
\coordinate (C) at (-3.0, 2.6);
\draw[gray!40, dotted, line width=0.8pt] (C) circle (0.48cm);
\node[dot, fill=blue!70!black] at (C) {};
\node[lbl, anchor=east] at ($(C)+(-0.55, 0.13)$) {$c$};
\node[sublbl, anchor=north east, align=right] at ($(C)+(-0.55,-0.13)$)
  {$c\in E,\;c\notin E'$\\[-1pt](isolated)};

% --- Legend ---
\begin{scope}[shift={(-2.6, -1.6)}]
  \fill[blue!12, draw=blue!50, line width=0.7pt] (0,0) rectangle (0.48,0.30);
  \node[lbl, anchor=west] at (0.62, 0.15) {$E$};

  \draw[orange!75!black, dashed, line width=1.2pt]
    (1.6, 0.15) -- (2.2, 0.15);
  \node[lbl, anchor=west] at (2.35, 0.15) {boundary of $E'$};

  \draw[gray!40, dotted, line width=0.8pt]
    (5.4, 0.15) -- (6.0, 0.15);
  \node[lbl, anchor=west] at (6.15, 0.15) {isolation ball around $c$};
\end{scope}

\end{tikzpicture}

\smallskip
{\small
  $a \in E \cap E'$: a non-isolated point of $E$, and a limit point of $E$.
  \quad
  $b \in E' \setminus E$: a limit point of $E$ not belonging to $E$.
  \quad
  $c \in E \setminus E'$: an isolated point of $E$, so $c \notin E'$.
}
\caption{The derived set $E' = \{x \in X : x \text{ is a limit point of } E\}$
relative to $E$.
The three annotated points illustrate the meaningful membership
combinations for points of~$X$.}
\label{fig:derived-set}
\end{figure}

% ---- Isolated Point ----
\begin{tcolorbox}[colback=propbox, colframe=propborder, arc=2pt,
  left=6pt, right=6pt, top=4pt, bottom=4pt,
  title={\small\textbf{Definition (Isolated Point)}},
  fonttitle=\small\bfseries]
\begin{definition}[Isolated Point]
\label{def:isolated-point}
Let $(X,d)$ be a metric space and let $E \subseteq X$.
A point $a \in E$ is called an \emph{isolated point} of $E$ if $a$ is not
a limit point of $E$, i.e.\ $a \notin E'$.
\end{definition}
\end{tcolorbox}

\begin{remark*}[Fully quantified form]
\[
\forall (X,d)\ \text{metric space},\
\forall E \subseteq X,\
\forall a \in E,
\]
\[
a\ \text{is isolated}
\iff
\exists r \in \mathbb{R}^+\ \text{such that}\
\forall y \in E,\ y \neq a \Rightarrow d(a,y) \geq r.
\]
\end{remark*}

\begin{remark*}[Isolated vs.\ limit point]
Isolated points are always \emph{elements} of $E$.
Limit points of $E$ need not belong to $E$: consider $E = (0,1) \subset \mathbb{R}$,
where $0$ and $1$ are limit points not in $E$.
Every point of $E$ is either isolated or a limit point of $E$ --- and both can
occur simultaneously in the same set (e.g.\ $E = \{0\} \cup (1,2)$).
\end{remark*}

% ---- Sequential Characterization ----
\begin{theorem}[\texorpdfstring{\hyperref[prf:seq-char-limit]{Sequential Characterization of Limit Points}}{Sequential Characterization of Limit Points}]
\label{thm:seq-char-limit}
Let $(X,d)$ be a metric space, $E \subseteq X$, and $x \in X$.
Then
\[
x \in E'
\iff
\exists (x_n) \subseteq E \text{ with all } x_n \neq x
\text{ and } x_n \to x.
\]
\end{theorem}

\begin{remark*}[Fully quantified form]
\[
\forall (X,d),\ \forall E \subseteq X,\ \forall x \in X,
\]
\[
x \in E'
\iff
\exists (x_n)_{n\in\mathbb{N}} \subseteq E
\text{ such that }
(\forall n,\ x_n \neq x)
\land
(\forall \varepsilon > 0,\ \exists N\ \forall n \ge N,\ d(x_n,x) < \varepsilon).
\]
\end{remark*}

% ---- Infinite Neighbourhood Characterization ----
\begin{corollary}[\texorpdfstring{\hyperref[prf:inf-nbhd]{Infinite Neighbourhood Characterization}}{Infinite Neighbourhood Characterization}]
\label{cor:inf-nbhd}
Let $(X,d)$ be a metric space and $E \subseteq X$.
Then
\[
x \in E'
\iff
\forall r > 0,\ B_r(x) \cap E
\text{ is infinite}.
\]
\end{corollary}

\begin{remark*}[Fully quantified form]
\[
x \in E'
\iff
\forall r > 0,\
\forall N \in \mathbb{N},\
\exists y \in B_r(x) \cap E
\text{ distinct from at least } N \text{ previously chosen points}.
\]
\end{remark*}

\begin{remark*}[Consequence]
This corollary rules out finite sets having limit points: if $E$ is finite
then every ball contains at most $|E|$ points of $E$, so no point can be a
limit point.
\end{remark*}

% ---- Derived Set is Closed ----
\begin{theorem}[\texorpdfstring{\hyperref[prf:derived-closed]{Derived Set is Closed}}{Derived Set is Closed}]
\label{thm:derived-closed}
Let $(X,d)$ be a metric space and $E \subseteq X$.
Then $E'$ is a closed subset of $X$.
\end{theorem}

\begin{remark*}[Fully quantified form]
\[
\forall (X,d),\ \forall E \subseteq X,\
E' \text{ is closed}
\iff
\forall (x_n) \subseteq E',\ x_n \to x
\Rightarrow x \in E'.
\]
\end{remark*}

% ---- Closed Sets and Derived Sets ----
\begin{theorem}[\texorpdfstring{\hyperref[prf:closed-derived]{Closed Sets and Derived Sets}}{Closed Sets and Derived Sets}]
\label{thm:closed-derived}
Let $(X,d)$ be a metric space and $E \subseteq X$.
Then
\[
E \text{ is closed}
\iff
E' \subseteq E.
\]
\end{theorem}

\begin{remark*}[Fully quantified form]
\[
\forall (X,d),\ \forall E \subseteq X,
\]
\[
\bigl(
\forall (x_n) \subseteq E,\ x_n \to x \Rightarrow x \in E
\bigr)
\iff
\bigl(
\forall x \in X,\ x \in E' \Rightarrow x \in E
\bigr).
\]
\end{remark*}

\begin{remark*}[Intuition]
Closedness means a set contains its own boundary crowding points.
Containment of the derived set is precisely this: every point $E$ accumulates
toward is already inside $E$.
\end{remark*}

% ---- Closure Formula ----
\begin{theorem}[\texorpdfstring{\hyperref[prf:closure-formula]{Closure Formula}}{Closure Formula}]
\label{thm:closure-formula}
Let $(X,d)$ be a metric space and $E \subseteq X$.
Then
\[
\overline{E} = E \cup E'.
\]
\end{theorem}

\begin{remark*}[Fully quantified form]
\[
\forall x \in X,\quad
x \in \overline{E}
\iff
\bigl(x \in E\bigr)
\lor
\bigl(x \in E'\bigr).
\]
\end{remark*}

\begin{remark*}[Consequence]
The closure adds exactly the missing limit points.
Since $E'$ is closed (Theorem~\ref{thm:derived-closed}),
$E \cup E'$ is the smallest closed set containing $E$,
which is precisely the definition of $\overline{E}$.
\end{remark*}

% ---- Perfect Set ----
\begin{tcolorbox}[colback=propbox, colframe=propborder, arc=2pt,
  left=6pt, right=6pt, top=4pt, bottom=4pt,
  title={\small\textbf{Definition (Perfect Set)}},
  fonttitle=\small\bfseries]
\begin{definition}[Perfect Set]
\label{def:perfect-set}
Let $(X,d)$ be a metric space.
A set $E \subseteq X$ is called \emph{perfect} if
\[
E \text{ is closed}
\quad \text{and} \quad
E = E'.
\]
\end{definition}
\end{tcolorbox}

\begin{remark*}[Intuition]
A perfect set has no isolated points (every point is a limit point) and no
missing limit points (closedness).
The Cantor set is the canonical example:
closed, nowhere dense, and every point is a limit point.
\end{remark*}


% =========================================================
% Derived Set E' and its Connection to Closure
% Source: volume-iii/analysis/metric-spaces/notes/limit-points/
%
% Purpose: Connect E' to closure explicitly with worked example
%          computing E', cl(E), int(E), ext(E), boundary(E)
%          for the concrete set E = (0,1) ∪ {2} in ℝ.
% Dependencies: def:limit-point, def:derived-set, def:closure
% =========================================================

\subsubsection{The Derived Set and Closure: How They Fit Together}
\label{sec:derived-set-closure-connection}

The \emph{derived set} $E'$ collects all limit points of $E$.
The \emph{closure} $\overline{E}$ is the smallest closed set containing $E$.
These are closely related, and the connection is made precise by
Theorem~\ref{thm:closure-formula}: $\overline{E} = E \cup E'$.

Before proving anything, let us build intuition with a concrete example
that we will compute from scratch.

\begin{example}[Computing $E'$, $\overline{E}$, $\partial E$, $\operatorname{int}(E)$,
$\operatorname{ext}(E)$ for $E = (0,1) \cup \{2\}$ in $\mathbb{R}$]
\label{ex:derived-set-full-computation}

Let $X = \mathbb{R}$ with the standard metric $d(x,y) = |x - y|$,
and let
\[
E = (0,1) \cup \{2\}.
\]

\medskip
\noindent\textbf{Step 1: Find the limit points $E'$.}

A point $x \in \mathbb{R}$ is a limit point of $E$ if every punctured
ball $B_r(x) \setminus \{x\}$ contains a point of $E$.

\begin{itemize}
  \item \textbf{If $x \in (0,1)$}: Take any $r > 0$.
        Then $(x - r, x + r) \setminus \{x\}$ contains points of $(0,1)$
        (since $(0,1)$ is an open interval with no isolated points).
        So every $x \in (0,1)$ is a limit point.

  \item \textbf{If $x = 0$}: Take any $r > 0$.
        Then $B_r(0) \setminus \{0\} = (-r, r) \setminus \{0\}$,
        which contains $(0, \min(r,1))$, a non-empty subset of $E$.
        So $0 \in E'$.

  \item \textbf{If $x = 1$}: Take any $r > 0$.
        Then $B_r(1) \setminus \{1\} = (1-r, 1+r) \setminus \{1\}$,
        which contains $(1-r, 1) \cap (0,1) \neq \emptyset$ for small $r$.
        So $1 \in E'$.

  \item \textbf{If $x = 2$}: Take $r = 1/2$.
        Then $B_{1/2}(2) \setminus \{2\} = (3/2, 5/2) \setminus \{2\}$.
        This interval contains no point of $E$:
        $(0,1) \cup \{2\}$ has no point in $(3/2, 5/2) \setminus \{2\}$.
        So $2 \notin E'$. \quad (The isolated point $2$ is not a limit point.)

  \item \textbf{If $x < 0$ or $x > 2$ or $x \in (1,2)$}: Similar reasoning shows these
        are not limit points (they are too far from $E$).
\end{itemize}

\noindent\textbf{Conclusion:} $E' = [0, 1]$.

\medskip
\noindent\textbf{Step 2: Closure $\overline{E} = E \cup E'$.}
\[
\overline{E} = (0,1) \cup \{2\} \cup [0,1] = [0,1] \cup \{2\}.
\]
This is the smallest closed set containing $E$.
Note: $[0,1]$ is closed (it contains its limit points $0$ and $1$),
and $\{2\}$ is closed (it has no limit points to miss).
Their union is closed.

\medskip
\noindent\textbf{Step 3: Interior $\operatorname{int}(E)$.}

A point $x \in E$ is interior if some $B_r(x) \subseteq E$.

\begin{itemize}
  \item For $x \in (0,1)$: take $r = \min(x, 1-x)/2 > 0$; then $B_r(x) \subseteq (0,1) \subseteq E$.
        So $(0,1) \subseteq \operatorname{int}(E)$.
  \item For $x = 2$: every $B_r(2)$ contains points outside $E$ (e.g.\ $2 + r/2$).
        So $2 \notin \operatorname{int}(E)$.
\end{itemize}

\noindent\textbf{Conclusion:} $\operatorname{int}(E) = (0,1)$.

\medskip
\noindent\textbf{Step 4: Boundary $\partial E$.}

$x \in \partial E$ iff every ball around $x$ meets both $E$ and $\mathbb{R} \setminus E$.

\begin{itemize}
  \item $x = 0$: every $B_r(0)$ contains points in $(0, r) \subseteq E$ and
        points in $(-r, 0) \subseteq \mathbb{R} \setminus E$. So $0 \in \partial E$.
  \item $x = 1$: every $B_r(1)$ contains $(1-r, 1) \subseteq E$ and $(1, 1+r) \subseteq \mathbb{R}\setminus E$.
        So $1 \in \partial E$.
  \item $x = 2$: every $B_r(2)$ contains $2 \in E$ and points in $(2-r, 2) \subseteq \mathbb{R} \setminus E$.
        So $2 \in \partial E$.
  \item $x \in (0,1)$: take $r = \min(x,1-x)/2$; then $B_r(x) \subseteq (0,1) \subseteq E$,
        so $B_r(x)$ misses $\mathbb{R} \setminus E$. So interior points are not boundary points.
\end{itemize}

\noindent\textbf{Conclusion:} $\partial E = \{0, 1, 2\}$.

\medskip
\noindent\textbf{Step 5: Exterior $\operatorname{ext}(E) = \operatorname{int}(\mathbb{R} \setminus E)$.}
\[
\operatorname{ext}(E) = (-\infty, 0) \cup (1, 2) \cup (2, \infty).
\]
(These are the open regions outside $E$.)

\medskip
\noindent\textbf{Summary table:}
\[
\begin{array}{ll}
E & = (0,1) \cup \{2\} \\
E' & = [0,1] \\
\overline{E} = E \cup E' & = [0,1] \cup \{2\} \\
\operatorname{int}(E) & = (0,1) \\
\partial E & = \{0, 1, 2\} \\
\operatorname{ext}(E) & = (-\infty,0) \cup (1,2) \cup (2,\infty)
\end{array}
\]

\noindent\textbf{Check: $\mathbb{R} = \operatorname{int}(E) \sqcup \partial E \sqcup \operatorname{ext}(E)$.}
\[
(0,1) \;\cup\; \{0,1,2\} \;\cup\; \bigl((-\infty,0)\cup(1,2)\cup(2,\infty)\bigr)
= \mathbb{R}. \checkmark
\]
\end{example}

\begin{remark*}[The isolated point $\{2\}$: not a limit point but still a boundary point]
A common confusion: the isolated point $2$ is \emph{not} in $E'$ (it has no
sequence in $E \setminus \{2\}$ converging to it), but it \emph{is} in $\partial E$
(every ball around $2$ hits $E$ because $2 \in E$, and hits $\mathbb{R} \setminus E$
because there are points near $2$ not in $E$).

Closure picks up both limit points and isolated points through $\overline{E} = E \cup E'$:
the isolated point $2$ comes from the $E$ part, not the $E'$ part.
\end{remark*}

\begin{remark*}[How $E'$ relates to $\overline{E}$ precisely]
Theorem~\ref{thm:closure-formula} gives the exact relationship:
\[
\overline{E} = E \cup E'.
\]
This says:
\begin{itemize}
  \item Every point already in $E$ is in $\overline{E}$ (trivially).
  \item Every limit point of $E$ is in $\overline{E}$, even if it is not in $E$.
  \item Nothing else is in $\overline{E}$.
\end{itemize}
In the example above, $0$ and $1$ are in $\overline{E}$ via $E'$
(they are limit points of $E$ not in $E$), while $2$ is in $\overline{E}$
via the $E$ part (it is in $E$ but not in $E'$).
\end{remark*}

